The previous proof had the right geometric idea; the gaps are in a few standard facts. Here is a fully explicit proof.

We work modulo \(n\) on the indices. For a nonzero vector \(u\) and a nonempty compact convex set \(K\subset \mathbb{R}^2\), write
\[
h_K(u):=\max_{x\in K}\langle u,x\rangle,
\qquad
F_K(u):=\{x\in K:\langle u,x\rangle=h_K(u)\}.
\]

## Preliminary facts

### 1. \(P+Q\) is a polygon
Since \(P=\operatorname{conv}(S)\) and \(Q=\operatorname{conv}(T)\), we have
\[
P+Q=\operatorname{conv}(S)+\operatorname{conv}(T)=\operatorname{conv}(S+T),
\]
where
\[
S+T:=\{s+t:s\in S,\ t\in T\}.
\]

Indeed, if
\[
p=\sum_{s\in S}\alpha_s s\in \operatorname{conv}(S),
\qquad
q=\sum_{t\in T}\beta_t t\in \operatorname{conv}(T),
\]
with \(\alpha_s,\beta_t\ge 0\) and \(\sum_s\alpha_s=\sum_t\beta_t=1\), then
\[
p+q=\sum_{s,t}\alpha_s\beta_t(s+t),
\]
so \(p+q\in \operatorname{conv}(S+T)\). Thus
\[
\operatorname{conv}(S)+\operatorname{conv}(T)\subseteq \operatorname{conv}(S+T).
\]
The reverse inclusion is immediate because \(S+T\subseteq P+Q\) and \(P+Q\) is convex.

Hence \(P+Q\) is the convex hull of finitely many points, i.e. a convex polygon. Also \(P+Q\) is \(2\)-dimensional, because for any \(q\in Q\),
\[
P+q\subseteq P+Q,
\]
and \(P\) is a genuine \(n\)-gon.

### 2. Support functions and exposed faces of Minkowski sums
For compact convex \(A,B\subset\mathbb{R}^2\) and \(u\neq 0\),
\[
h_{A+B}(u)=h_A(u)+h_B(u).
\]
This is immediate:
\[
h_{A+B}(u)
=\max_{a\in A,\ b\in B}\langle u,a+b\rangle
=\max_{a,b}\bigl(\langle u,a\rangle+\langle u,b\rangle\bigr)
=h_A(u)+h_B(u).
\]

From this one gets
\[
F_{A+B}(u)=F_A(u)+F_B(u).
\]
Proof: if \(a\in F_A(u)\) and \(b\in F_B(u)\), then
\[
\langle u,a+b\rangle=h_A(u)+h_B(u)=h_{A+B}(u),
\]
so \(a+b\in F_{A+B}(u)\). Thus
\[
F_A(u)+F_B(u)\subseteq F_{A+B}(u).
\]
Conversely, if \(z\in F_{A+B}(u)\), write \(z=a+b\) with \(a\in A\), \(b\in B\). Then
\[
h_A(u)+h_B(u)=h_{A+B}(u)=\langle u,z\rangle
=\langle u,a\rangle+\langle u,b\rangle\le h_A(u)+h_B(u),
\]
so equality holds in both inequalities, hence
\[
a\in F_A(u),\qquad b\in F_B(u).
\]
Therefore \(z\in F_A(u)+F_B(u)\).

### 3. Unique maximizers on \(T\) are also unique maximizers on \(Q=\operatorname{conv}(T)\)
If \(t^\ast\in T\) is the unique maximizer of \(\langle u,\cdot\rangle\) on \(T\), then
\[
F_Q(u)=\{t^\ast\}.
\]
Indeed, for any
\[
q=\sum_{t\in T}\alpha_t t\in Q
\qquad
(\alpha_t\ge 0,\ \sum_t\alpha_t=1),
\]
we have
\[
\langle u,q\rangle=\sum_{t\in T}\alpha_t\langle u,t\rangle\le \langle u,t^\ast\rangle,
\]
with equality only if every \(t\) with \(\alpha_t>0\) also maximizes \(\langle u,\cdot\rangle\) on \(T\). By uniqueness, that forces \(q=t^\ast\).

### 4. A singleton exposed face is a vertex
Let \(K\) be a convex polygon. If \(F_K(u)=\{x\}\), then \(x\) is a vertex of \(K\).

Proof: if
\[
x=\lambda a+(1-\lambda)b
\qquad
(0<\lambda<1,\ a,b\in K),
\]
then
\[
h_K(u)=\langle u,x\rangle
=\lambda \langle u,a\rangle+(1-\lambda)\langle u,b\rangle
\le \lambda h_K(u)+(1-\lambda)h_K(u)=h_K(u).
\]
So equality holds throughout, hence \(a,b\in F_K(u)=\{x\}\), and therefore \(a=b=x\). Thus \(x\) is extreme, i.e. a vertex.

### 5. Every vertex of a convex polygon has an open cone of exposing directions
Let \(K\) be a \(2\)-dimensional convex polygon and let \(y\) be a vertex. Let \(E^-\) and \(E^+\) be the two edges of \(K\) meeting at \(y\), and let \(n^-,n^+\) be outer normals to those edges. Then
\[
\langle n^-,z-y\rangle\le 0,\qquad \langle n^+,z-y\rangle\le 0
\quad\text{for all }z\in K,
\]
with equality in the first exactly on \(E^-\), and equality in the second exactly on \(E^+\).

Hence for every
\[
u=\alpha n^-+\beta n^+ \qquad (\alpha,\beta>0),
\]
and every \(z\in K\setminus\{y\}\), at least one of the two inequalities is strict, so
\[
\langle u,z-y\rangle<0.
\]
Therefore
\[
F_K(u)=\{y\}.
\]
So the set
\[
D_y:=\{\alpha n^-+\beta n^+:\alpha,\beta>0\}
\]
is a nonempty open cone of directions uniquely exposing \(y\).

---

## Part (2)

Fix \(i\). Let
\[
E_i=[s_i,s_{i+1}],\qquad a_i:=s_{i+1}-s_i,
\]
and let \(u_i\) be an outer normal to \(E_i\). Set
\[
G_i:=F_Q(u_i).
\]

We will prove
\[
G_i=[t_i,t_{i+1}],
\]
and then deduce the direction statement.

### Step 1: small perturbations of \(u_i\) enter the adjacent normal cones
Because \(u_i\) is an outer normal to the edge \(E_i\),
\[
F_P(u_i)=E_i=[s_i,s_{i+1}],
\]
so
\[
\langle u_i,s_i\rangle=\langle u_i,s_{i+1}\rangle=h_P(u_i),
\]
and for every \(j\neq i,i+1\),
\[
\langle u_i,s_j\rangle<h_P(u_i).
\]
Define
\[
\delta_j:=h_P(u_i)-\langle u_i,s_j\rangle>0
\qquad (j\neq i,i+1),
\]
and let
\[
\delta:=\min_{j\neq i,i+1}\delta_j>0.
\]
Also define
\[
M:=\max_{j\neq i,i+1}
\max\Bigl\{
|\langle a_i,s_i-s_j\rangle|,
|\langle a_i,s_{i+1}-s_j\rangle|
\Bigr\}.
\]
Choose \(0<\varepsilon<\delta/(M+1)\).

Then for \(j\neq i,i+1\),
\[
\langle u_i-\varepsilon a_i,\ s_i-s_j\rangle
=
\delta_j-\varepsilon\langle a_i,s_i-s_j\rangle
\ge \delta-\varepsilon M>0,
\]
and
\[
\langle u_i-\varepsilon a_i,\ s_i-s_{i+1}\rangle
=
-\varepsilon\langle a_i,a_i\rangle
\cdot(-1)
=
\varepsilon\|a_i\|^2>0.
\]
So \(s_i\) is the unique maximizer of \(\langle u_i-\varepsilon a_i,\cdot\rangle\) on \(P\), i.e.
\[
u_i-\varepsilon a_i\in C_i.
\]

Similarly, for \(j\neq i,i+1\),
\[
\langle u_i+\varepsilon a_i,\ s_{i+1}-s_j\rangle
=
\delta_j+\varepsilon\langle a_i,s_{i+1}-s_j\rangle
\ge \delta-\varepsilon M>0,
\]
and
\[
\langle u_i+\varepsilon a_i,\ s_{i+1}-s_i\rangle
=
\varepsilon\|a_i\|^2>0.
\]
Hence
\[
u_i+\varepsilon a_i\in C_{i+1}.
\]

So for all sufficiently small \(\varepsilon>0\),
\[
u_i-\varepsilon a_i\in C_i,
\qquad
u_i+\varepsilon a_i\in C_{i+1}.
\]

### Step 2: \(t_i,t_{i+1}\in G_i\)
Take any sequence \(\varepsilon_m\downarrow 0\) such that
\[
u_i-\varepsilon_m a_i\in C_i.
\]
Since \(t_i\) is the unique maximizer on \(T\) for every direction in \(C_i\), for every \(t\in T\),
\[
\langle u_i-\varepsilon_m a_i,\ t_i-t\rangle\ge 0.
\]
Passing to the limit \(m\to\infty\) gives
\[
\langle u_i,\ t_i-t\rangle\ge 0
\qquad \text{for all }t\in T.
\]
Thus \(t_i\) maximizes \(\langle u_i,\cdot\rangle\) on \(T\), hence on \(Q=\operatorname{conv}(T)\). Therefore
\[
t_i\in G_i.
\]

Exactly the same argument, using
\[
u_i+\varepsilon_m a_i\in C_{i+1},
\]
shows that
\[
t_{i+1}\in G_i.
\]

### Step 3: \(G_i\) is exactly the segment \([t_i,t_{i+1}]\)
Since \(G_i=F_Q(u_i)\), it lies in the support line
\[
L_i:=\{q\in\mathbb{R}^2:\langle u_i,q\rangle=h_Q(u_i)\}.
\]
Because \(u_i\perp a_i\), the line \(L_i\) is parallel to \(a_i\). Hence \(G_i\) is a compact convex subset of a line parallel to \(a_i\), so \(G_i\) is either a point or a closed segment parallel to \(a_i\).

Now fix a sufficiently small \(\varepsilon>0\) with
\[
u_i-\varepsilon a_i\in C_i,\qquad u_i+\varepsilon a_i\in C_{i+1}.
\]

For any \(q\in G_i\), since \(\langle u_i,q\rangle=h_Q(u_i)\),
\[
\langle u_i-\varepsilon a_i,q\rangle
=
h_Q(u_i)-\varepsilon\langle a_i,q\rangle.
\]
So **among points of \(G_i\)**, maximizing \(\langle u_i-\varepsilon a_i,\cdot\rangle\) is equivalent to minimizing \(\langle a_i,\cdot\rangle\).

But \(u_i-\varepsilon a_i\in C_i\), so by Subproblem \(2\) and Preliminary Fact \(3\),
\[
F_Q(u_i-\varepsilon a_i)=\{t_i\}.
\]
Since \(t_i\in G_i\), it follows that \(t_i\) is the unique minimizer of \(\langle a_i,\cdot\rangle\) on \(G_i\).

Similarly, for \(q\in G_i\),
\[
\langle u_i+\varepsilon a_i,q\rangle
=
h_Q(u_i)+\varepsilon\langle a_i,q\rangle,
\]
so **among points of \(G_i\)**, maximizing \(\langle u_i+\varepsilon a_i,\cdot\rangle\) is equivalent to maximizing \(\langle a_i,\cdot\rangle\). Since
\[
F_Q(u_i+\varepsilon a_i)=\{t_{i+1}\},
\]
the point \(t_{i+1}\) is the unique maximizer of \(\langle a_i,\cdot\rangle\) on \(G_i\).

Therefore the two endpoints of the segment \(G_i\) are exactly \(t_i\) and \(t_{i+1}\). Hence
\[
G_i=[t_i,t_{i+1}].
\]

Since faces do not change when the normal is multiplied by a positive scalar, this holds for **any** outer normal to \(E_i\).

### Step 4: direction of \(t_{i+1}-t_i\)
Because \(G_i=[t_i,t_{i+1}]\) lies on a line parallel to \(a_i\), there exists \(\lambda_i\in\mathbb{R}\) such that
\[
t_{i+1}-t_i=\lambda_i a_i=\lambda_i(s_{i+1}-s_i).
\]
Also \(t_i\) is the minimum and \(t_{i+1}\) the maximum of \(\langle a_i,\cdot\rangle\) on \(G_i\), so
\[
\langle a_i,t_{i+1}-t_i\rangle\ge 0.
\]
Substituting \(t_{i+1}-t_i=\lambda_i a_i\) gives
\[
\lambda_i\|a_i\|^2\ge 0.
\]
Since \(a_i\neq 0\), we get
\[
\lambda_i\ge 0.
\]

This proves Part \((2)\).

---

## Part (1)

Define
\[
x_i:=s_i+t_i.
\]

### Step 1: each \(x_i\) is a vertex of \(P+Q\)
Fix \(i\) and let \(u\in C_i\). Then
\[
F_P(u)=\{s_i\},
\]
because \(u\) lies in the open normal cone of \(P\) at \(s_i\). By Subproblem \(2\) and Preliminary Fact \(3\),
\[
F_Q(u)=\{t_i\}.
\]
Therefore, by the face formula,
\[
F_{P+Q}(u)=F_P(u)+F_Q(u)=\{s_i+t_i\}=\{x_i\}.
\]
By Preliminary Fact \(4\), \(x_i\) is a vertex of \(P+Q\).

Also, by Subproblems \(1\)–\(2\), the points \(\ell_i=x_i\) are the \(n\) support points in the extremal case, so they are pairwise distinct.

### Step 2: there are no other vertices
Let \(y\) be any vertex of \(P+Q\). By Preliminary Fact \(5\), there exists a nonempty open cone \(D_y\) such that
\[
F_{P+Q}(u)=\{y\}
\qquad\text{for all }u\in D_y.
\]

For the polygon \(P\), every nonzero direction lies either in one of the open vertex cones \(C_1,\dots,C_n\), or on one of the \(n\) boundary rays determined by the edge normals \(u_1,\dots,u_n\). The union of those boundary rays has empty interior, so the open cone \(D_y\) meets some \(C_i\).

Take \(u\in D_y\cap C_i\). Then
\[
\{y\}=F_{P+Q}(u)=\{x_i\},
\]
hence
\[
y=x_i.
\]
So every vertex of \(P+Q\) is one of the \(x_i\).

Since each \(x_i\) is a vertex and there are no others, the vertices of \(P+Q\) are exactly
\[
x_1,\dots,x_n.
\]

### Step 3: cyclic order
For the edge normal \(u_i\), Part \((2)\) and the face formula give
\[
F_{P+Q}(u_i)
=
F_P(u_i)+F_Q(u_i)
=
[s_i,s_{i+1}]+[t_i,t_{i+1}].
\]
Write
\[
t_{i+1}-t_i=\lambda_i a_i,\qquad \lambda_i\ge 0.
\]
Then
\[
[s_i,s_{i+1}]
=
\{s_i+\alpha a_i:0\le \alpha\le 1\},
\]
and
\[
[t_i,t_{i+1}]
=
\{t_i+\beta \lambda_i a_i:0\le \beta\le 1\}.
\]
Hence
\[
[s_i,s_{i+1}]+[t_i,t_{i+1}]
=
\{x_i+(\alpha+\lambda_i\beta)a_i:0\le \alpha,\beta\le 1\}
=
[x_i,x_{i+1}],
\]
because \(x_{i+1}=x_i+(1+\lambda_i)a_i\).

So for each \(i\),
\[
F_{P+Q}(u_i)=[x_i,x_{i+1}],
\]
which is a boundary edge of \(P+Q\). Therefore \(x_i\) and \(x_{i+1}\) are adjacent vertices for every \(i\). Since the vertices are exactly the pairwise distinct points \(x_1,\dots,x_n\), this means they occur on the boundary in the cyclic order
\[
x_1,x_2,\dots,x_n.
\]

That proves Part \((1)\).

---

## Conclusion

We have shown:

\[
\boxed{\text{The vertices of }P+Q\text{ are exactly }x_i=s_i+t_i\text{ in cyclic order.}}
\]

And for each \(i\), if \(u_i\) is any outer normal to the edge \([s_i,s_{i+1}]\) of \(P\), then
\[
\boxed{F_Q(u_i)=[t_i,t_{i+1}]}
\]
and
\[
\boxed{t_{i+1}-t_i=\lambda_i(s_{i+1}-s_i)\quad\text{for some }\lambda_i\ge 0.}
\]

So Part \((2)\) also says that the edge of \(P+Q\) corresponding to the edge \([s_i,s_{i+1}]\) of \(P\) is exactly
\[
[x_i,x_{i+1}],
\]
parallel to \([s_i,s_{i+1}]\).